\documentclass[12pt,a4paper]{article}
\usepackage{amssymb,amsmath}
\usepackage[T2A]{fontenc}
\usepackage[utf8]{inputenc}
\usepackage[russian]{babel}

\usepackage{epsfig,graphicx}

% \input ris.tex
% \addrisoption{\risleftfalse}

\usepackage{pict2e}
\usepackage{tikz}

\DeclareRobustCommand{\No}{\ifmmode{\nfss@text{\textnumero}}\else\textnumero\fi}

\catcode`\@=11
\def\kratno{\mathrel{\smash{\lower.5ex\hbox{$\,\vdots\,$}}}}
\def\nekratno{\mathrel{\mathpalette\c@ncel\kratno}}
\def\c@ncel#1#2{\m@th\ooalign{$\hfil#1\mkern1mu/\hfil$\crcr$#1#2$}}
\def\divis{\smash{\lower.1ex\hbox{\,\hbox{.}\kern-.22em\lower-.6ex\hbox{.}
\kern-.52em\lower-1.2ex\hbox{.}\,}} }

\hyphenation{чис-ло чис-ла чис-лу чис-лом чис-ле чис-лам чис-ла-ми чис-лах
че-ты-рех-уголь-ник че-ты-рех-уголь-ни-ка пря-мо-уголь-ник тре-уголь-ни-ка
тре-уголь-ни-ков вне-впи-сан-ной вне-впи-сан-ных}

\sloppy
\pagestyle{empty}

\def\q#1.{{\bf #1. }}

\def\btable#1#2#3\endbtable{{%
\offinterlineskip\def\\{\cr\noalign{\hrule}}\halign
  {\strut\vrule\hbox to #1{\hfil#2\hfil}\vrule width 1.5 pt&&\hbox to #1{\hfil#2\hfil}\vrule\cr\noalign{\hrule}#3}}}

\begin{document}

\centerline{\sc САНКТ-ПЕТЕРБУРГСКАЯ}
\centerline{\sc ОЛИМПИАДА ШКОЛЬНИКОВ 2025 года ПО МАТЕМАТИКЕ}
\centerline{\sc I тур. \ 6 класс.}
\medskip
\hrule
\bigskip

\q1. Галя, Толя и Валя записали в тетрадях одну и ту же строку из 10 различных букв. Каждый стёр несколько букв в своей тетради. У Гали осталось слово БАСНЯ, у Толи --- ДЕНЬ. Могло ли у Вали остаться слово ДЯТЕЛ? Если да~--- напишите исходную последовательность букв, если нет --- объясните, почему не могло.

\q2. Докажите, что никакую прямоугольную таблицу не удастся покрасить в три цвета так,
чтобы в каждой строке было 12~красных и 16 синих клеток (а остальные зелёные), и в то же время в~каждом столбце~--- 24 синих и 20 зеленых клеток (а остальные красные).

\q3. Арсений написал 10 троек, дописал справа четверку, потом справа ещё 9 троек, дописал справа четверку, снова 9 троек, снова четверку, и так несколько раз, чередуя группы из 9 троек и четверки, а~в~конце числа после очередных 9 троек дописал единицу. Докажите, что получилось составное число. 

\q4. Мастер закрепляет на стене прямоугольный металлический лист, после чего при помощи станка вырезает из него прямоугольные детали. У листа и всех вырезаемых деталей одна сторона вертикальна, другая горизонтальна. Все вырезаемые станком детали одинаковые, то есть у них равны вертикальные стороны и равны горизонтальные стороны. (Например, если станок вырезает детали $2\,\text{м}\times 3$\,м, то вырезать деталь $3\,\text{м}\times 2$\,м он не может.)

При помощи этого станка из листа $10\,\text{м}\times 10\,\text{м}$ как-то вырезали 10~деталей. Докажите, что из любого прямоугольного металлического листа, стороны которого не меньше 10\,м, а площадь равна 220 квадратных метров, мастеру удастся вырезать 11 деталей. (Длины сторон деталей и~второго листа не обязательно целые.)


\clearpage

\centerline{\sc САНКТ-ПЕТЕРБУРГСКАЯ}
\centerline{\sc ОЛИМПИАДА ШКОЛЬНИКОВ 2025 года ПО МАТЕМАТИКЕ}
\centerline{\sc I тур. \ 7 класс.}
\medskip
\hrule
\bigskip

\q1. Из $80$ белых кубиков сложили параллелепипед $4\times 4\times 5$, после чего 
покрасили его снаружи в красный цвет.
Можно ли теперь из тех же 80 кубиков сложить параллелепипед $8\times 10\times 1$ так, чтобы одна из его граней размером $8\times 10$ оказалась полностью красной?

\q2. Существует ли натуральное число, оканчивающееся на 34, у~которого делителей, оканчивающихся на 8, больше, чем делителей, оканчивающихся на 9?

\q3. На олимпиаду пришло 557 детей.  Их как-то рассадили по трём аудиториям (в каждой аудитории есть хотя бы по одному ребенку). В каждой аудитории подсчитали, какой процент от находящихся в ней детей составляют девочки. Сумма трёх полученных чисел оказалась равна 280. Найдите наименьшее возможное количество девочек среди этих 557 детей.

\q4. В каждой клетке квадратного поля $10\times 10$ стоит печенег или хазар. Печенеги всегда говорят правду, а хазары каждое число увеличивают на 1 (например, если хазар хочет сказать <<четыре>>, он произносит <<пять>>). Каждого спросили, сколько печенегов среди его соседей по стороне, сложили все ответы и получили сумму 292. Затем каждого спросили, сколько хазар среди его соседей по стороне, сложили все ответы и получили сумму 140. 

\hskip10pt
{\bf а)} Сколько всего хазар стоит на этом поле? 

\hskip10pt
{\bf б)} Сколько имеется способов расставить печенегов и хазар на этом поле так, чтобы получились такие суммы?


\clearpage

\centerline{\sc САНКТ-ПЕТЕРБУРГСКАЯ}
\centerline{\sc ОЛИМПИАДА ШКОЛЬНИКОВ 2025 года ПО МАТЕМАТИКЕ}
\centerline{\sc I тур. \ 8 класс.}
\medskip
\hrule
\bigskip

\q1. Найдите три таких различных трёхзначных числа $a$, $b$ и $c$, что каждое из них равно сумме какого-то делителя $a$, какого-то делителя $b$ и какого-то делителя $c$. Достаточно привести один пример такой тройки чисел.

\q2. Найдите все тройки положительных чисел $x, y, z$, удовлетворяющих уравнениям
$$
\frac x2=\frac1y+\frac 1z,\qquad  \frac y2=\frac1z+\frac 1x,\qquad \frac z2=\frac1x+\frac1y.
$$

\q3. На равных сторонах $AB$ и $BC$ равнобедренного треугольника $ABC$ отмечены точки $D$ и $E$ соответственно так, что $AE = CD$. Точка $H$ на отрезке $BE$ — основание высоты, опущенной из точки~$A$ на сторону $BC$. Найдите длину отрезка $EH$, если известно, что $CE = 2$, $AD = 4$.

\q4. В классе учится чётное число учеников, все они разного роста. На каждом уроке они садятся за парты по двое. На уроках алгебры, геометрии и искусственного интеллекта оказалось, что для каждой парты сумма ростов сидящих за ней учеников равна 3 м, 3,3 м или 3,5 м. Докажите, что какие-то двое сидели за одной партой хотя бы на двух из этих трёх уроков. 

\q5. Даны натуральные числа $m$ и $n$, $m<n$. В десятичной записи дроби $m/n$ после запятой подряд встретились цифры $2, 0, 2, 4$ (именно в таком порядке). Докажите, что существует правильная дробь со знаменателем $n$ (и натуральным числителем), в записи которой после запятой найдутся подряд цифры 1 и 2 именно в~таком порядке.


\clearpage

\centerline{\sc САНКТ-ПЕТЕРБУРГСКАЯ}
\centerline{\sc ОЛИМПИАДА ШКОЛЬНИКОВ 2025 года ПО МАТЕМАТИКЕ}
\centerline{\sc I тур. \ 9 класс.}
\medskip
\hrule
\bigskip

\q1. У Пети, Васи и Толи в тетрадях было записано одно и то же десятизначное число. Каждый стёр несколько цифр в своей тетради. У Пети получилось число 12436, у Васи~--- 3578. Могло ли у Толи получиться число 9510?

\q2. Найдите все тройки ненулевых чисел $x, y, z$, каждое из которых в 3 раза меньше суммы чисел, обратных к двум другим.

\q3. В трапеции $ABCD$ диагональ $BD$ равна основанию $AD$, а~также $\angle A = 2 \angle D$ и $AB = 2BC$. Докажите, что $\angle ACD = 90^\circ$.

\q4. Клетчатый прямоугольник периметра $p$ удалось разрезать на 100 клетчатых прямоугольников, никакие два из которых не равны. У каждого из них есть сторона длины 2. Найдите наименьшее возможное значение $p$. Прямоугольники, которые отличаются поворотом, считаются одинаковыми.

\q5. В бесконечной возрастающей последовательности $a_1, a_2, \ldots$ натуральных чисел любые два соседних числа отличаются не более чем на миллион. Верно ли, что можно выбрать миллион членов этой последовательности так, чтобы их наибольший общий делитель был больше миллиона?


\clearpage

\centerline{\sc САНКТ-ПЕТЕРБУРГСКАЯ}
\centerline{\sc ОЛИМПИАДА ШКОЛЬНИКОВ 2025 года ПО МАТЕМАТИКЕ}
\centerline{\sc I тур. \ 10 класс.}
\medskip
\hrule
\bigskip

\q1. Решите уравнение 
$$
\rule[-5.2pt]{0.5pt}{17.4pt}\hspace{3pt}
\rule[-4.4pt]{0.5pt}{15.8pt}\hspace{3pt}
\dots
\rule[-3.6pt]{0.5pt}{14.2pt}\hspace{3pt}
\rule[-2.8pt]{0.5pt}{12.6pt}\hspace{3pt}
\rule[-2pt]{0.5pt}{11pt}\hspace{1.5pt}2x
\hspace{1.5pt}\rule[-2pt]{0.5pt}{11pt}
-x
\hspace{1.5pt}\rule[-2.8pt]{0.5pt}{12.6pt}
-x
\hspace{1.5pt}\rule[-3.6pt]{0.5pt}{14.2pt}
-\dots
\hspace{1.5pt}\rule[-4.4pt]{0.5pt}{15.8pt}
-x
\hspace{1.5pt}\rule[-5.2pt]{0.5pt}{17.4pt}
=x^2-1.
$$
В левой части знак <<минус>> фигурирует 98 раз.

\q2. Дан треугольник $ABC$, в котором $\angle A = 70^\circ$. На
стороне  $AB$ отмечена точка $X$, на стороне $BC$ --- точка~$Y$, а на
стороне $AC$ --- точка $Z$, причем $AB = BY$, $CY = CZ$ и $AZ =AX$.
Найдите угол $XYZ$.

\q3. На доске написано некоторое натуральное число $N$. Петя разделил его с остатком на 4441, записал в тетрадь остаток, а полученное неполное частное разделил на 81 и снова записал остаток. Вася разделил $N$ с остатком на 81, записал к себе в тетрадь остаток, полученное неполное частное разделил на 4441 и снова записал остаток в тетрадь. Сумма двух остатков, записанных  Петей, оказалась не равна сумме двух Васиных остатков. На какое наименьшее число могли отличаться друг от друга эти суммы?

\q4. Сумма целых чисел $a$ и $b$ не равна 1. Известно, что число $n^2-2an-b$ не делится на $a+b-1$ ни при каком целом~$n$. Докажите, что квадратный трехчлен $x^2-2bx-a$ не имеет целых корней.

\q5. Какое наименьшее количество клеток можно отметить в квадрате $110\times
110$ так, чтобы в любом прямоугольнике $11\times 12$ (и в любом прямоугольнике
$12\times 11$) была хотя бы одна отмеченная клетка?


\clearpage

\centerline{\sc САНКТ-ПЕТЕРБУРГСКАЯ}
\centerline{\sc ОЛИМПИАДА ШКОЛЬНИКОВ 2025 года ПО МАТЕМАТИКЕ}
\centerline{\sc I тур. \ 11 класс.}
\medskip
\hrule
\bigskip

\q1. Каждое из 170 натуральных чисел равно сумме чисел, обратных к 169 остальным. Найдите эти числа. (Укажите все варианты и докажите, что других нет.)

\q2. В остроугольном треугольнике $ABC$ проведены высота $BH$, биссектриса $BL$ и медиана $BM$ (точки $A,H,L,M,C$ расположены на прямой именно в таком порядке). Оказалось, что $MH=1$ и~радиус описанной окружности треугольника равен~2. Найдите угол $\angle ALB$.

\q3. На доске написано чётное количество различных вещественных чисел. Дима, Витя и Саша разбили числа на пары (каждый по-своему, разные мальчики не могли выбрать одну и ту же пару), после чего перемножили числа в парах. Каждое из полученных произведений оказалось равно 2023, 2024 или 2025. Докажите, что кто-то из мальчиков ошибся.

\q4. Сумма целых чисел $a$ и $b$ не равна 1. Известно, что число $n^2-2an-b$ не делится на $a+b-1$ ни при каком целом~$n$. Докажите, что квадратный трехчлен $x^2-2bx-a$ не имеет целых корней.

\q5. Куб со стороной $2^{1\,000\,000}$ разбит на единичные кубики. В каждом кубике написана цифра 0, 1 или 2. Назовем строку из цифр 0, 1 и 2 \emph{хорошей}, если ее можно получить, начав с некоторого кубика и переходя на каждом шагу к соседнему (по грани) кубику. Например, строка 0102 хорошая, если можно из кубика с цифрой 0 попасть в кубик с цифрой 1, из него~--- в кубик с нулем (возможно, начальный), а~из него~--- в кубик с двойкой.

Оказалось, что любая строка длины не более 10\,000~--- хорошая. Один из кубиков удалили, из-за чего некоторые строки перестали быть хорошими.  Докажите, что осталось не менее $2\cdot 6^{1012}-2^{2024}$ хороших строк длины 2024.

\end{document}
